\chapter{Limits and Open Problems}
\label{ch:limits}
% src: formal/twin-containment/README.md (not covered); docs/AUTONOMY_PLAN.md section 6;
% docs/STRATA_TAXONOMY.md; paper/aies2027/kit/CALIBRATION.md; owner decisions of 2026-10-04.

The structure keeps a misreading from becoming permission. It does not keep the robot from
misreading, and in our runs the misreadings were common. This chapter states what the results rest
on, what the proofs leave out and which choices in the home model belong to people rather than to
engineering.

\section{What the results rest on}

The evaluation uses one simulated home, one set of rules that we wrote and scenarios that we wrote
or that a separate agent wrote from our protocol. The simulated world reports itself truthfully to
the robot's perception. A camera that misses Margaret on the floor, a microphone that mishears her
or a vision model that invents a person would add failures this evaluation cannot show. The
held-out set is blind by procedure. A separate agent wrote it and only its hashes were committed
before the system was frozen, but no third party wrote or kept it.

The development runs before the harness fix measured the test harness as much as the robot. In
those runs 44 of the 63 check-ins recorded as unanswered followed an answer Margaret had given. The
class accuracy of every run before that fix is confounded, including some of its successes, since a
lost answer also pushed real emergencies up the ladder. The containment and privacy results of
those runs are not confounded, because they judge the robot's own actions against the rulings it
had.

One harness defect is still open. The runner opens each scenario's window by counting the robot's
records once the world is reset, and a decision that lands in between falls outside the window. In
dev13 that happened once, in d36, and the grader flagged an authorized emergency call as a breach.
A grader that reads the decision service's ledger, rather than the runner's window, would not
miss it. Until that change, a flagged breach is checked by hand against the ledger.

\section{What the proofs leave out}

The Lean model treats every model output as arbitrary and proves that none reaches authority. It
says nothing about sensors that lie below their signatures. A device whose key is stolen, or a
physical event staged in front of two working sensors, produces valid witnesses. It says nothing
about who writes the home model and its bars, and a careless author can make any action
unguarded. Two learned models reach authority, and the proofs cover neither. The monitoring centre
reads Margaret's own words with a model, so a misread request for help can send emergency services.
The robot's own messages cannot reach that decision, but hers can. The camera's person detector is a
learned vision model whose output is one witness. A detector that sees a fall that did not happen
adds a false witness, and when the centre cannot be reached that one witness is enough to call
emergency services. In the twin the decision service attests the detector's output itself and does
not check that signature, so the camera witness is as trustworthy as the service that runs it. Static tests tie the model to the
code by requiring the real home model to have the model's shape. They do not prove that the Python
implements the model.

\section{Choices that belong to people}

Several rules in the home model are policy, and the owner set them on 2026-10-04 after the
development runs. The robot never lifts a person who has fallen, except to move her away from a
danger in the home, because a hidden injury can be made worse. The dispatcher sends police only on
stated facts of a crime, such as an assault, a forced entry, a demand for her money or a theft,
and never for a stranger merely present. The robot is a mandatory reporter. An articulable
suspicion of a crime against Margaret obliges it to report to the monitoring centre before
anything else, and no refusal of hers waives the report. Each choice trades her autonomy against
her safety or against a third party's interest, and people who deploy such a robot should make
them deliberately.

Two thresholds are of the same kind. Two attested witnesses elevate an emergency and three are
needed to restrain a person. A minute below the bar ends a right. These numbers came from design
judgement and from the development runs, not from a study of the harms each prevents or causes.

\section{Open problems}

The world can change between a decision and the act, as when an animal leaves or responders enter.
Every graded run in Chapter~\ref{ch:evaluation} judged an action at the record that chose it, so
those runs show permission at the decision, not at the act. The robot now checks again at the act
(Section~\ref{sec:home-act}). That check is built and tested but was not in any graded run, and the
held-out run will be the first to measure it. It closes the gap inside the robot, where one ledger
orders every event. When the monitoring centre and the dispatcher keep ledgers of their own,
ordering their events against the robot's will need vector clocks or a similar causal order.

The other kinds of machine that can enter a home, a peer robot, a delivery drone or an unknown
robot, currently pass through the visitor rules and the network-message rules. A household with
several robots needs strata of its own for them.

Finally, the preregistered study of an internal monitor for the robot's classifier
(Chapter~\ref{ch:ieip}) asks whether a failure of representation equivariance in the classifier's
hidden states predicts its classification errors. If it does, the robot could check in rather than
act on a reading its own classifier is unsure of. The preregistration states its prior that the
monitor will fail, as its consumer-relative variant did, and the outcome is reported either way.
